\subsection{Radon spaces}

DEFINITION 2. --- Let T be a topological space. One says that T is a Radon (resp. strongly Radon) space if T is Hausdorff and if every function defined on the Borel tribe \(\mathfrak{B}(\mathrm{T})\) of T, with values in \(\overline{R}_{+}\) , that is countably additive and bounded (resp. locally bounded) is inner regular.

For example, we shall see later on (Prop. 3) that every Polish space is strongly Radon. In particular, every locally compact space with a countable base is strongly Radon.

There exist Radon spaces that are not strongly Radon.

PROPOSITION 1. --- Every Lindelöf \(^{(1)}\) Radon space is strongly Radon.

Let T be a Lindelöf space that is Radon, and let I be a set function on the tribe \(\mathfrak{B}(T)\) that is positive, countably additive and locally bounded. The open sets V such that \(\mathrm{I}(\mathrm{V}) < +\infty\) form a covering of T, from which one can extract a countable covering \((\mathrm{V}_{n})_{n \in \mathbb{N}}\) . Set \(G_{n} = V_{0} \cup V_{1} \cup \cdots \cup V_{n}\) for every \(n \in N\) ; set \(H_{0} = G_{0}\) and \(H_{n} = G_{n} - G_{n-1}\) for \(n \geqslant 1\) ; finally, denote by \(I_{n}\) the set function \(A \mapsto \mathrm{I}(A \cap H_{n})\) on \(\mathfrak{B}(T)\) , which is obviously countably additive and bounded. Since the sets \(H_{n}\) form a partition of T, we have \(I = \sum_{n} I_{n}\) . The space T being Radon, for every \(n \in N\) there exists a bounded measure \(\mu_{n}\) on T such that \(\mu_{n}^{\bullet}(\mathrm{A}) = I_{n}(\mathrm{A})\) for all \(A \in \mathfrak{B}(T)\) ; therefore also \(\sum_{n} \mu_{n}^{\bullet}(\mathrm{A}) = I(\mathrm{A})\) . Since I is locally bounded, the family \((\mu_{n})\) is summable (§1, No.~7, Prop. 7); if \(\mu\) denotes \(\sum_{n} \mu_{n}\) , we have \(\mu^{\bullet}(\mathrm{A}) = I(\mathrm{A})\) for all \(A \in \mathfrak{B}(T)\) , and it follows that I is inner regular. In other words, T is strongly Radon.

Recall that a subset A of a topological space T is said to be universally measurable if A is \(\mu\) -measurable for every measure \(\mu\) on T. This is equivalent to saying that A is \(\mu\) -measurable for every measure \(\mu\) on T with compact support ( \(\S1\) , No.~8, Prop. 9).

PROPOSITION 2. --- Let X be a topological space and T a subspace of X.

\begin{enumerate}
\def\labelenumi{\alph{enumi})}
\tightlist
\item
  Suppose that \(\mathbf{T}\) is a Radon space. Then, for every function I defined on \(\mathfrak{B}(\mathbf{X})\) that is positive, countably additive and bounded, one has
\end{enumerate}

\[
\sup_{\substack{K\text{compact}\\ K\subset T}}I(K) = \inf_{\substack{B\in \mathfrak{B}(X)\\ B\supset T}}I(B).\tag{6}
\]

Moreover, \(\mathbf{T}\) is universally measurable in \(\mathbf{X}\) .

\begin{enumerate}
\def\labelenumi{\alph{enumi})}
\setcounter{enumi}{1}
\tightlist
\item
  Conversely, suppose that X is a Radon space and that T is universally measurable in X; then T is a Radon space.
\end{enumerate}

Let us prove \(a\) ). We denote by \(\alpha\) the right side of (6); for every \(n \in \mathbf{N}\) , there exists a set \(C_n \in \mathfrak{B}(X)\) containing \(T\) , such that \(I(C_n) \leqslant \alpha + 2^{-n}\) . Setting \(C = \bigcap_{n} C_n\) , we then have \(T \subset C\) , \(I(C) = \alpha\) . If \(A \in \mathfrak{B}(T)\) , let us choose a Borel set \(B\) of \(X\) such that \(A = B \cap T\) (GT, IX, §6, No.~3) and set \(J(A) = I(B \cap C)\) . This number does not depend on the choice of \(B\) , for if \(B'\) is a second Borel set in \(X\) such that \(A = B' \cap T\) , then \(B \cap C\) and \(B' \cap C\) differ only by a Borel set \(M\) contained in \(C - T\) , and \(I(M) = 0\) by the construction of \(C\) . Clearly \(J(K) = I(K)\) for every compact set \(K \subset T\) . Let \((A_n)\) be a sequence of Borel sets of \(T\) , pairwise disjoint, and, for each \(n\) , let \(B_n\) be a Borel set of \(X\) such that \(B_n \cap T = A_n\) . Replacing \(B_n\) by \(B_n - \left( \bigcup_{k < n} B_k \right)\) if necessary, we can suppose that the sets \(B_n\) are pairwise disjoint. Set \(A = \bigcup_{n} A_n\) and \(B = \bigcup_{n} B_n\) ; then

\[
\mathrm{J} (\mathrm{A}) = \mathrm{I} (\mathrm{B} \cap \mathrm{C}) = \sum_ {n} \mathrm{I} (\mathrm{B} _ {n} \cap \mathrm{C}) = \sum_ {n} \mathrm{J} (\mathrm{A} _ {n});
\]

J is thus a countably additive and bounded function on \(\mathfrak{B}(\mathrm{T})\) . Since T is by hypothesis a Radon space, there exists a bounded measure \(\mu\) on T such that \(\mathrm{J}(\mathrm{A}) = \mu^{\bullet}(\mathrm{A})\) for every \(\mathrm{A} \in \mathfrak{B}(\mathrm{T})\) ; consequently

\[
\alpha = \mathrm{J} (\mathrm{T}) = \mu^ {\bullet} (\mathrm{T}) = \sup _ {\mathrm{K}} \mu^ {\bullet} (\mathrm{K}) = \sup _ {\mathrm{K}} \mathrm{J} (\mathrm{K}),
\]

by the definition of \(\mu^{\bullet}\) . Formula (6) is thus established.

Let us now show that T is universally measurable. Let \(\lambda\) be a bounded measure on X; the preceding argument may be applied to the set function \(I: A \mapsto \lambda^{\bullet}(A)\) on \(\mathfrak{B}(X)\) , thus there exists a sequence \((K_{n})\) of compact subsets of T such that (with the above notations)

\[
\sup _ {n} \lambda^ {\bullet} (\mathrm{K} _ {n}) = \mathrm{J} (\mathrm{T}) = \lambda^ {\bullet} (\mathrm{C}).
\]

Set \(K' = \bigcup_{n \in \mathbf{N}} K_n\) ; \(K'\) is Borel in \(X\) , \(K' \subset T \subset C\) , \(\lambda^\bullet(K') = \lambda^\bullet(C)\) , therefore these three sets differ only by \(\lambda\) -negligible sets, and so \(T\) is \(\lambda\) -measurable. This completes the proof of \(a\) ).

Let us pass to \(b\) ). Suppose that \(X\) is a Radon space, and that \(T\) is universally measurable in \(X\) . Let \(I\) be a positive function on \(\mathfrak{B}(T)\) that is countably additive and bounded; the function \(A \mapsto I(A \cap T)\) on \(\mathfrak{B}(X)\) is then positive, countably additive and bounded, therefore there exists a bounded measure \(\nu\) on \(X\) such that \(I(A \cap T) = \nu^{\bullet}(A)\) for all \(A \in \mathfrak{B}(X)\) . Now, \(T\) is \(\nu\) -measurable; the preceding relation shows that \(\nu^{\bullet}(K) = 0\) for every compact subset \(K\) of \(X\) that is disjoint from \(T\) , therefore \(\nu\) is concentrated on \(T\) . Consequently, for every Borel set \(A\) of \(X\) , we have \(I(A \cap T) = \nu^{\bullet}(A \cap T) = \mu^{\bullet}(A \cap T)\) , where \(\mu\) is the measure induced by \(\nu\) on \(T\) . Finally, it follows that \(I(B) = \mu^{\bullet}(B)\) for every set \(B \in \mathfrak{B}(T)\) (GT, IX, §6, No.~3, Remark 2), and \(I\) is indeed inner regular.

COROLLARY. --- If X is a Radon space, then every Borel subset T of X is Radon.

For, \(\mathbf{T}\) is universally measurable in \(\mathbf{X}\) .

PROPOSITION 3. --- Every Souslin space (in particular, every Polish or Lusin space) is strongly Radon.

Let T be a Souslin space; since T is a Lindelöf space (TG, IX, Appendix I, Cor. of Prop. 1), \(^{(2)}\) it suffices to show that T is Radon (Prop. 1). Let I be a function defined on \(\mathfrak{B}(T)\) , positive, countably additive and bounded. We extend I to \(\mathfrak{P}(T)\) by setting, for every subset A of T,

\[
\operatorname{I}(\mathbf{A}) = \inf_{\substack{\mathbf{B}\in \mathfrak{B}(\mathbf{T})\\ \mathbf{B}\supset \mathbf{A}}}\operatorname{I}(\mathbf{B}).
\]

Let us show that this extension is a capacity on T (GT, IX, §6, No.~9). It is clear that the relation \(\mathbf{A} \subset \mathbf{A}'\) implies \(\mathrm{I}(\mathbf{A}) \leqslant \mathrm{I}(\mathbf{A}')\) . Let \((\mathbf{A}_n)\) be an increasing sequence of subsets of T, and let \(\mathbf{A} = \bigcup_{n} \mathbf{A}_n\) . The set of Borel sets that contain \(\mathbf{A}_n\) being stable for countable intersections, there exists for each \(n\) a Borel set \(\mathbf{B}_n\) such that \(\mathbf{A}_n \subset \mathbf{B}_n\) and \(\mathrm{I}(\mathbf{A}_n) = \mathrm{I}(\mathbf{B}_n)\) (cf.~the proof of Prop. 2). Set \(\mathbf{C}_n = \bigcap_{p \geqslant n} \mathbf{B}_p\) ; \(\mathbf{C}_n\) is Borel, and \(\mathbf{A}_n \subset \mathbf{C}_n \subset \mathbf{B}_n\) , therefore \(\mathrm{I}(\mathbf{A}_n) = \mathrm{I}(\mathbf{C}_n)\) . On the other hand, the sequence \((\mathbf{C}_n)\) is increasing. Let \(\mathbf{C} = \bigcup_{n} \mathbf{C}_n\) : the relation \(\mathbf{A} \subset \mathbf{C}\) implies that

\[
\mathrm{I} (\mathrm{A}) \leqslant \mathrm{I} (\mathrm{C}) = \lim _ {n} \mathrm{I} (\mathrm{C} _ {n}) = \lim _ {n} \mathrm{I} (\mathrm{A} _ {n}),
\]

whence the equality \(\mathrm{I}(\mathrm{A})=\lim_{n}\mathrm{I}(\mathrm{A}_{n})\) is immediate. Consequently, I is a capacity.

If \((\mathbf{H}_{n})\) is a decreasing sequence of closed sets in T, obviously \(\mathrm{I}\left(\bigcap_{n}\mathrm{H}_{n}\right)=\inf_{n}\mathrm{I}(\mathrm{H}_{n})\) . It follows that every Souslin subset F of T is capacitable for I (TG, IX, §6, No.~10, Prop. 15). In particular, every Borel set A of T is capacitable (loc. cit., §6, No.~3, Prop. 10). \(^{(3)}\) In other words,

\[
\mathrm{I} (\mathrm{A}) = \sup _ {\mathrm{K}} \mathrm{I} (\mathrm{K}),
\]

where K runs over the set of compact sets contained in A; we have proved that I is inner regular.

Remark. --- Let X be a Lusin space (in particular, any Polish space), and f a bijective continuous mapping of X onto a (Lusin) regular space Y. One knows (TG, IX, §6, No.~7, Prop. 14) that the mapping \(B \mapsto f^{-1}(B)\) is a bijection of the Borel tribe of Y onto the Borel tribe of X. The spaces X and Y are Lusin, hence strongly Radon (Prop. 3). It follows immediately that the mapping \(\mu \mapsto f(\mu)\) is a bijection of the set of bounded measures on X onto the set of bounded measures on Y. \(^{(4)}\)

